\documentclass[a4paper,12pt]{article}
\usepackage[utf8]{inputenc}
\usepackage[table,xcdraw,x11names]{xcolor} % Para los colores en las celdas de la tabla y las transparencias
\usepackage{graphicx}
\usepackage{transparent} % Paquete para aplicar transparencia a imágene
\usepackage{fancyhdr}
\usepackage{geometry}
\usepackage{lastpage} % Para obtener el número total de páginas
\usepackage[absolute,overlay]{textpos} % Paquete para posiciones absolutas
\usepackage{atbegshi} % Paquete para controlar elementos en cada página
\usepackage{eso-pic} % Paquete para añadir elementos en el fondo de la página
\definecolor{verdeoscuro}{rgb}{0.0, 0.2, 0.0}
\usepackage{longtable} % Para tablas de varias páginas
\usepackage{array} % Para usar el entorno array dentro de celdas
\usepackage{multirow}
\usepackage{multicol}
\usepackage[hidelinks]{hyperref}                                     %Para añadir enlaces.
\geometry{a4paper, top=2cm, left=2cm, right=2cm, bottom=4.8cm}
\setlength{\headheight}{70pt} % Ajusta la altura del encabezado según sea necesario
\setlength{\footskip}{14.5pt} % Ajusta la distancia del pie de página al texto
% Eliminar la indentación al principio de cada párrafo
\setlength{\parindent}{0pt}
% Establecer la distancia estándar entre párrafos
\setlength{\parskip}{1.5em}

% Configuración del encabezado y pie de página
\pagestyle{fancy}
\fancyhf{}

% Encabezado personalizado
\fancyhead[L]{
    \begin{minipage}[b]{1.5cm}
        \includegraphics[height=50pt]{imagenes/logo_cole_transp.png} % Logo o imagen1 en el encabezado
    \end{minipage}
    \hspace{0.2cm} % Espacio entre logo y texto
    \begin{minipage}[b]{6cm}
        \footnotesize{\textbf{CEIP MAR MEDITERRÁNEO} \\
        Consejería de Desarrollo Educativo \\
        y Formación Profesional \\
        \vspace{0.01cm}}
    \end{minipage}
    \hspace{0.1cm} % Espacio entre el texto y la imagen europa.png
    \begin{minipage}[b]{1.5cm}
        \includegraphics[height=50pt]{imagenes/europa.png} % Imagen2 a la derecha del texto
        \vspace{0.01cm}
    \end{minipage}
    \hspace{2.5cm} % Espacio entre las dos imágenes
    \begin{minipage}[b]{2.5cm}
        \includegraphics[height=30pt]{imagenes/junta.png} % Imagen3 a la derecha de imagen2
        \vspace{-0.3cm} % Ajusta la posición vertical del número de página
        \makebox[4cm][r]{\footnotesize{Página \thepage\ de \pageref{LastPage}}} % Número de página alineado a la derecha
        \vspace{0.01cm}
    \end{minipage}
}

% Configuración del texto en el pie de página usando 'fancyhdr'
\fancyfoot[L]{
    \begin{minipage}[t]{7cm} % Minipage para controlar el área del texto
        \raggedright % Alinea el texto a la izquierda
        \rule{7cm}{0.1pt} \\
        \textcolor{verdeoscuro}{\footnotesize Calle José Morales Abad, 9. 04007 Almería} \\
        \textcolor{verdeoscuro}{Tlf.: 950 15 62 03} \\
        \textcolor{verdeoscuro}{04005326.edu@juntadeandalucia.es}
    \end{minipage}
}

% Añadir la imagen escudo.png como fondo en todas las páginas en la esquina inferior derecha
\AddToShipoutPictureBG{%
    \AtPageLowerLeft{%
        \put(\LenToUnit{\dimexpr\paperwidth-2.8cm\relax},\LenToUnit{0.2cm}){%
            \transparent{0.2}\includegraphics[height=100pt]{imagenes/escudo.png} % Ajusta el tamaño de la imagen con transparencia
        }
    }
}

\renewcommand{\contentsname}{Índice}  % Cambia 'Contents' por 'Índice'

\begin{document}

\begin{center}
    \section*{DIARIO DE CLASE}
\end{center}

\tableofcontents

\newpage

\sloppy % Para partir las palabras

\section*{Repaso de cuarto}

\subsection{20240918 - Repaso de cuarto.}

\begin{enumerate}
    \item Las tablas de multiplicar.
    \item Mi primer programa en \textit{Python}.
\end{enumerate}

\subsection{20240919 - Repaso de cuarto.}

\begin{enumerate}
    \item Las tablas de multiplicar.
    \item Mi primer programa en \textit{Python}.
    \item Los números.
    \begin{itemize}
        \item Los números de seis cifras.
        \item El valor de las cifras de un número.
        \item Comparación y ordenación de números.
        \item Aproximación de números.
        \item Los números ordinales.
        \item Los números romanos.
    \end{itemize}
    \item La suma y la resta.
    \begin{itemize}
        \item La suma y sus propiedades.
        \item La resta. La prueba de la resta.
        \item El paréntesis.
        \item Cálculo estimado.
    \end{itemize}
    \item La multiplicación.
    \begin{itemize}
        \item La multiplicación y sus términos.
        \item Las tablas de multiplicar.
        \item Propiedades de la multiplicación.
        \item Propiedad distributiva.
        \item La multiplicación por una cifra.
        \item La multiplicación por varias cifras.
        \item La multiplicación por la unidad seguida de ceros.
    \end{itemize}
\end{enumerate}

\subsection{20240920 - Repaso de cuarto.}

\begin{enumerate}
    \item Las tablas de multiplicar.
    \item Mi primer programa en \textit{Python}.
    \item Estadística y azar.
    \begin{itemize}
        \item Datos cualitativos y datos cuantitativos.
        \item Tablas de registro de datos.
        \item Pictogramas.
        \item Gráficos de barras.
        \item Gráficos de líneas.
        \item Experiencias aleatorias y sucesos.
    \end{itemize}
    \item La división.
    \begin{itemize}
        \item La división como reparto.
        \item La división como partición.
        \item La prueba de la división.
        \item La división con ceros en el cociente.
        \item La división entre la unidad seguida de ceros.
        \item Operaciones combinadas.
    \end{itemize}
    \item Las fracciones.
    \begin{itemize}
        \item Las fracciones y sus términos.
        \item Medios, tercios y cuartos.
        \item Nombramos a las fraciones.
        \item Comparación de fracciones.
    \end{itemize}
\end{enumerate}

\subsection{20240925 - Repaso de cuarto}

\begin{enumerate}
    \item Prueba de nivel de las tablas de multiplicar.
    \item Los número decimales y el dinero.
    \begin{itemize}
        \item Fracción decimal y número decimal.
        \item Uso de monedas y billetes.
        \item Comparación y ordenación de decimales.
        \item Aproximación de números decimales.
        \item Suma y resta de números decimales.
        \item Operaciones con euros y céntimos.
    \end{itemize}
    \item Geometría y orientación en el plano.
    \begin{itemize}
        \item Rectas, semirectas y segmentos.
        \item Rectas paralelas y secantes.
        \item Los ángulos.
        \item Los ángulos según su abertura.
        \item El ángulo como giro.
        \item Traslación.
        \item Simetría.
        \item Simulación en el plano: coordenadas.
        \item Orientación espacial.
    \end{itemize}
\end{enumerate}

\subsection{20240926 - Repaso de cuarto}

\begin{enumerate}
    \item La medida del tiempo.
    \begin{itemize}
        \item El año: el calendario.
        \item Lustros, décadas, siglos y milenios.
        \item Días, horas, minutos y segundos.
        \item Relojes analógicos y digitales.
    \end{itemize}
    \item La medida de la longitud.
    \begin{itemize}
        \item El metro y sus múltiplos.
        \item Los submúltiplos del metro.
        \item Expresiones complejas e incomplejas.
        \item Operaciones con medidas de longitud.
    \end{itemize}
\end{enumerate}

\subsection{20240927 - Repaso de cuarto}

\begin{enumerate}
    \item La medida de la capacidad y de la masa.
    \begin{itemize}
        \item La medida de la capacidad.
        \item La medida de la masa.
        \item Expresiones complejas e incomplejas.
        \item Operaciones con capacidades y masas.
    \end{itemize}
    \item Figuras planas y cuerpos geométricos.
    \begin{itemize}
        \item Los polígonos.
        \item Clasificación de polígonos.
        \item Clasificación de triángulos.
        \item Clasificación de cuadriláteros.
        \item Introducción al área.
        \item Circunferencia y círculo.
        \item Poliedros. Prismas y pirámides.
        \item Cuerpos redondos.
    \end{itemize}
\end{enumerate}

\newpage

\section{Primer Trimestre}

\subsection{20241002 - Sesión 01 - Libro S1 y S2}

\begin{itemize}
    \item Conozo el valor posicional de cada cifra.
    \begin{itemize}
        \item Actividades 1, 2 y 3 página 3.
        \item Opcional: Juego Reagrupo números, página 4.
    \end{itemize}
    \item Estimo y mido.
    \begin{itemize}
        \item Lectura: \textit{En busca de lombrices y diamantes I}. (Página 5).
        \item Respondemos a las preguntas de la página 6.
        \item Actividades 1 y 2 página 7.
    \end{itemize}
    \item Programo en \textit{Python}.
    \begin{itemize}
        \item Uso el comando \textbf{\textit{print}} para imprimir texto por pantalla.
    \end{itemize}
\end{itemize}

\subsection{20241003 - Sesión 02 - Libro S3 y S4}

\begin{itemize}
    \item Revisión del trabajo de la sesión anterior.
    \item Represento diagramas estadísticos.
    \begin{itemize}
        \item Actividades 1 y 2 página 8; 3 página 9.
    \end{itemize}
    \item Utilizo paréntesis.
    \begin{itemize}
        \item Analizamos la toma de decisiones, página 10.
        \item Actividades 1, 2, 3 y 4 página 11.
        \item Opcional: Juego Cubo 21.
    \end{itemize}
    \item Programo en \textit{Python}.
    \begin{itemize}
        \item Uso de variables para guardar datos.
    \end{itemize}
\end{itemize}

\subsection{20241004 - Sesión 03 - Libro S5 y S6}

\begin{itemize}
    \item Revisión del trabajo de la sesión anterior.
    \item Calculo distancias.
    \begin{itemize}
        \item Actividades 1 página 13; 2 y 3 página 14.
    \end{itemize}
    \item Multiplico centenas por centenas.
    \begin{itemize}
        \item Actividades 1, 2 y 3 página 16.
    \end{itemize}
    \item Programo en \textit{Python}.
    \begin{itemize}
        \item Uso el comando \textbf{\textit{print}} para imprimir por pantalla texto mezclado con el contenido de variables.
    \end{itemize}
\end{itemize}

\subsection{20241009 - Sesión 04 - Libro S7 y S8}

\begin{itemize}
    \item Revisión del trabajo de la sesión anterior.
    \item Interpreto datos.
    \begin{itemize}
        \item Actividades 1 página 17 y 2 página 18.
    \end{itemize}
    \item Multiplico dos números naturales con más de tres cifras.
    \begin{itemize}
        \item Lectura: \textit{En busca de lombrices y diamantes II.} (Página 19).
        \item Respondemos a las preguntas de la página 20.
        \item Actividades 1 y 2 página 21.
    \end{itemize}
    \item Programo con \textit{Python.}
    \begin{itemize}
        \item ¿Puedo hacer que \textit{Python} me escriba cualquier tabla de multiplicar I?
    \end{itemize}
\end{itemize}

\subsection{20241010 - Sesión 05 - Libro S9 y S10}

\begin{itemize}
    \item Revisión del trabajo de la sesión anterior.
    \item Aplico la multiplicación.
    \begin{itemize}
        \item Actividades 1, 2 y 3 página 22.
        \item Opcional: Multiplicar con cuatro cubos, página 23.
    \end{itemize}
    \item Mido figuras rectangulares.
    \begin{itemize}
        \item Actividades 1, 2 y 3 página 25.
    \end{itemize}
    \item Programo con \textit{Python.}
    \begin{itemize}
        \item ¿Puedo hacer que \textit{Python} me escriba cualquier tabla de multiplicar II?
    \end{itemize}
\end{itemize}

\subsection{20241016 - Sesión 06 - Libro S11 y S12}

\begin{itemize}
    \item Revisión del trabajo de la sesión anterior.
    \item Interpreto restos.
    \begin{itemize}
        \item Lectura de la página 26.
        \item Actividad 1 página 27.
    \end{itemize}
    \item Lectura: \textit{En busca de lombrices-diamante III}. (páginas 28 y 29).
    \item Respondemos a las preguntas sobre la lectura de la página 29.
    \item Practico la división.
    \begin{itemize}
        \item Actividades 1, 2 y 3 página 30.
        \item Opcional: Cociente menor. (Página 31).
    \end{itemize}
    \item ¿Puedo hacer que \textit{Python} me escriba cualquier tabla de multiplicar III?
\end{itemize}

\subsection{20241017 - Sesión 07 - Libro S13 y S14}

\begin{itemize}
    \item Revisión del trabajo de la sesión anterior.
    \item Relaciono el dinero y los números decimales.
    \begin{itemize}
        \item Actividades 1 y 2 página 32; 3 y 4 página 33.
        \item Opcional: Reagrupo decimales. (Página 34).
    \end{itemize}
    \item Sumo, resto y comparo números decimales.
    \begin{itemize}
        \item Actividades 1, 2, 3 y 4 página 35; 5, 6, 7 y 8 página 36.
    \end{itemize}
    \item ¿Puedo hacer que \textit{Python} me escriba cualquier tabla de multiplicar IV?
\end{itemize}

\subsection{20241018 - Sesión 08 - Libro S15 y S16}

\begin{itemize}
    \item Lectura: \textit{En busca de lombrices-diamante IV}. (Páginas 37 y 38).
    \item Respondemos a las preguntas sobre la lectura de la página 38.
    \item Multiplico y divido decimales por potencias de 10.
    \begin{itemize}
        \item Actividades 1 y 2 página 39.
    \end{itemize}
    \item Aplico las aproximaciones.
    \begin{itemize}
        \item Actividades 1, 2, 3, 4 y 5 página 40.
    \end{itemize}
    \item ¿Puedo hacer que \textit{Python} me escriba cualquier tabla de multiplicar V?
\end{itemize}

\subsection{20241023 - Sesión 09 - Libro S13 y S14}

\begin{itemize}
    \item Revisión del trabajo de la sesión anterior.
    \item Repaso: Operaciones con decimales.
    \item Completamos las actividades que faltan por hacer de las páginas 32, 33, 35 y 36.
    \item Repasamos el bucle \textbf{\textit{for}} en \textit{Python}.
\end{itemize}

\subsection{20241024 - Sesión 10 - Libro S15 y S16}

\begin{itemize}
    \item Revisión del trabajo de la sesión anterior.
    \item Repaso: Divisiones por potencias de 10.
    \item Repaso: Aplicamos las aproximaciones.
    \item Completamos las actividades que faltan por hacer de las páginas 39 y 40.
    \item Repasamos el bucle \textbf{\textit{for}} en \textit{Python}.
\end{itemize}

\subsection{20241025 - Sesión 11 - Libro S17 y S18}

\begin{itemize}
    \item Revisión del trabajo de la sesión anterior.
    \item Lectura: \textit{En busca de lombrices-diamante V. (Páginas 41 y 42).}
    \item Respondemos a las preguntas sobre la lectura de la página 42.
    \item Relaciono unidades métricas.
    \begin{itemize}
        \item Actividades 1 página 43; 2, 3, 4 y 5 página 44.
    \end{itemize}
    \item Repaso los elementos de los polígonos.
    \begin{itemize}
        \item Actividades 1 y 2 página 45.
    \end{itemize}
\end{itemize}

\subsection{20241031 - Sesión 12 - Libro S19 y S20}

\begin{itemize}
    \item Revisión del trabajo de la sesión anterior.
    \item .
\end{itemize}
\end{document}